\section{Slide options}

% Switches work in the preamble (for the whole deck) or between frames (for
% the frames after them):
%   \TitleBarOn / \TitleBarOff          blue bar, or underlined black title
%   \SetTitleHeight{1.1cm}              height of the title area
%   \SectionNameOn / \SectionNameOff    section name at the top right
%   \FooterOn / \FooterNumberOnly / \FooterOff
%   \SetBarMargin{0.4cm}                side margin of title bar and footer
\TitleBarOff
\FooterNumberOnly
\begin{frame}{Slide options}{\cmd{TitleBarOff} and \cmd{FooterNumberOnly}}
  \begin{itemize}
    \item This frame follows \cmd{TitleBarOff} and \cmd{FooterNumberOnly}
    \item \cmd{TitleBarOn} and \cmd{FooterOn} switch back
    \item \cmd{SetTitleHeight\{1.1cm\}} changes the title area
    \item \cmd{showdividers} outlines rows and columns while editing
  \end{itemize}
\end{frame}
\TitleBarOn
\FooterOn

% Speaker notes: shown only in the presenter view of the HTML deck (press S)
% and not in the PDF. \note[item]{...} adds a bullet.
\begin{frame}{Speaker notes}
  \begin{itemize}
    \item Press \texttt{S} in the HTML deck for the presenter view
    \item It shows the notes, the next slide and a timer
    \item \texttt{F} full screen, \texttt{P} play/pause animations, \texttt{H} help,
          a number + \texttt{Enter} goes to that slide
  \end{itemize}
  \note{What to say on this slide.}
  \note[item]{A talking point}
  \note[item]{Another talking point}
\end{frame}
